\documentclass[a4paper,12pt]{article}
\usepackage[utf8]{inputenc}
\usepackage{hyperref}
\usepackage{geometry}
\usepackage{listings}
\usepackage{xcolor}
\usepackage{array}
\usepackage{booktabs}
\geometry{margin=1in}

\lstset{
    basicstyle=\ttfamily\small,
    keywordstyle=\color{blue},
    commentstyle=\color{green!50!black},
    stringstyle=\color{red},
    numbers=left,
    numberstyle=\tiny,
    frame=single,
    breaklines=true
}

\title{mighf System Specification \\ \large Complete Architecture, ISA, and Porting Guide v2.0}
\author{Miguel "Kuyo" / SufremOak}
\date{\today}

\begin{document}

\maketitle

\begin{abstract}
This document provides a comprehensive specification of the mighf system architecture, including detailed instruction set documentation, binary format specification, memory model, and complete guidelines for porting to new platforms. Based on the reference implementation mhfsoft.c, this specification ensures consistent behavior across all mighf implementations.
\end{abstract}

\tableofcontents
\newpage

\section{Overview}
The mighf system is a minimalist, 8-bit register / 16-bit address CPU architecture designed for educational use, embedded systems, and easy implementation across diverse hardware platforms. It features a simple instruction set, straightforward execution model, and well-defined binary format.

\section{Architecture Components}

\subsection{Register Set}
\begin{itemize}
    \item \textbf{R0--R7}: 8 general-purpose 8-bit registers for computation and data storage
    \item \textbf{PC} (Program Counter): 16-bit register holding the address of the next instruction
    \item \textbf{Memory}: 64KB (65536 bytes) flat address space, byte-addressable
\end{itemize}

\textbf{Important Note}: Unlike the original specification, the reference implementation uses 8-bit registers, not 16-bit. This design choice prioritizes simplicity and memory efficiency.

\subsection{Memory Model}
\begin{itemize}
    \item \textbf{Address Space}: 16-bit flat addressing (0x0000 - 0xFFFF)
    \item \textbf{Byte Addressable}: All memory operations work on byte granularity
    \item \textbf{Memory-Mapped I/O}: Address 0xFF00 is reserved for character output
    \item \textbf{Program Storage}: Code, data, and symbols can be loaded at arbitrary addresses
    \item \textbf{No Stack Pointer}: Current implementation does not include automatic stack management
\end{itemize}

\section{Binary Format Specification}

\subsection{MHF Binary Format (.mhfb)}
All mighf programs use the MHF binary format with the following header structure:

\begin{lstlisting}[language=C, caption=MHF Header Structure]
typedef struct {
    char signature[4];      // "MHF\0"
    uint16_t version;       // Format version
    uint16_t flags;         // Reserved flags
    uint32_t head_offset;   // Header data offset
    uint32_t code_offset;   // Code section offset (entry point)
    uint32_t data_offset;   // Data section offset
    uint32_t symbol_offset; // Symbol table offset
    uint32_t file_size;     // Total file size
} __attribute__((packed)) mhf_header_t;
\end{lstlisting}

\textbf{Loading Process}:
\begin{enumerate}
    \item Verify signature matches "MHF" (first 3 bytes)
    \item Read entire file content into memory starting at \texttt{head\_offset}
    \item Set PC to \texttt{code\_offset} as the entry point
\end{enumerate}

\section{Instruction Set Architecture}

\subsection{Instruction Format}
All instructions follow these patterns:
\begin{itemize}
    \item \textbf{Single Byte}: \texttt{[OPCODE]}
    \item \textbf{With Register}: \texttt{[OPCODE] [REG]}
    \item \textbf{With Immediate}: \texttt{[OPCODE] [REG] [IMM8]}
    \item \textbf{With Address}: \texttt{[OPCODE] [REG] [ADDR\_LOW] [ADDR\_HIGH]}
\end{itemize}

\subsection{Core Instruction Set}

\begin{table}[h]
\centering
\begin{tabular}{|l|l|l|p{6cm}|}
\hline
\textbf{Opcode} & \textbf{Mnemonic} & \textbf{Format} & \textbf{Description} \\
\hline
0x00 & NOP & - & No operation \\
\hline
0x01 & HALT & - & Stop CPU execution \\
\hline
0x10 & LOAD & reg, imm8 & Load 8-bit immediate into register \\
\hline
0x11 & STORE & reg, addr16 & Store register value to memory address \\
\hline
0x20 & ADD & reg1, reg2 & Add reg2 to reg1, store result in reg1 \\
\hline
\end{tabular}
\caption{mighf Core Instruction Set}
\end{table}

\subsection{Detailed Instruction Descriptions}

\subsubsection{0x00 - NOP}
\textbf{Format}: \texttt{0x00}\\
\textbf{Operation}: No operation, PC increments by 1\\
\textbf{Flags}: None affected

\subsubsection{0x01 - HALT}
\textbf{Format}: \texttt{0x01}\\
\textbf{Operation}: Stops CPU execution, sets running flag to 0\\
\textbf{Flags}: None affected\\
\textbf{Implementation Note}: Should output halt message for debugging

\subsubsection{0x10 - LOAD reg, imm8}
\textbf{Format}: \texttt{0x10 [REG] [IMM8]}\\
\textbf{Operation}: \texttt{R[reg] = imm8}\\
\textbf{Validation}: Register index must be 0-7\\
\textbf{Example}: \texttt{0x10 0x02 0x42} loads 66 into R2

\subsubsection{0x11 - STORE reg, addr16}
\textbf{Format}: \texttt{0x11 [REG] [ADDR\_LOW] [ADDR\_HIGH]}\\
\textbf{Operation}: \texttt{MEMORY[addr16] = R[reg]}\\
\textbf{Address Calculation}: \texttt{addr16 = (ADDR\_HIGH << 8) | ADDR\_LOW}\\
\textbf{Special Behavior}: Writing to 0xFF00 outputs character to console\\
\textbf{Example}: \texttt{0x11 0x00 0x00 0xFF} stores R0 to address 0xFF00 (console output)

\subsubsection{0x20 - ADD reg1, reg2}
\textbf{Format}: \texttt{0x20 [REG1] [REG2]}\\
\textbf{Operation}: \texttt{R[reg1] = R[reg1] + R[reg2]}\\
\textbf{Validation}: Both register indices must be 0-7\\
\textbf{Overflow}: 8-bit arithmetic with wrap-around\\
\textbf{Example}: \texttt{0x20 0x01 0x02} adds R2 to R1, stores result in R1

\section{Execution Model}

\subsection{CPU State Structure}
\begin{lstlisting}[language=C, caption=Reference CPU State]
typedef struct {
    uint8_t regs[8];        // 8-bit general purpose registers
    uint16_t pc;            // 16-bit program counter
    uint8_t memory[65536];  // 64KB memory space
    int running;            // Execution state flag
} MighfCPU;
\end{lstlisting}

\subsection{Execution Loop}
\begin{lstlisting}[language=C, caption=Standard Execution Loop]
void run_cpu(MighfCPU *cpu, uint16_t entry_point) {
    cpu->pc = entry_point;
    cpu->running = 1;
    
    while (cpu->running) {
        uint8_t opcode = cpu->memory[cpu->pc++];
        
        // Decode and execute instruction
        switch (opcode) {
            // Instruction implementations...
        }
    }
}
\end{lstlisting}

\subsection{Error Handling}
Implementations must handle:
\begin{itemize}
    \item \textbf{Invalid Opcodes}: Stop execution, report error with PC location
    \item \textbf{Register Bounds}: Validate register indices (0-7)
    \item \textbf{Memory Bounds}: Validate memory addresses (0x0000-0xFFFF)
    \item \textbf{File Format}: Verify MHF signature and structure
\end{itemize}

\section{Memory-Mapped I/O}

\subsection{Console Output (0xFF00)}
Writing any value to address 0xFF00 outputs the value as a character:
\begin{itemize}
    \item \textbf{Implementation}: Use platform-appropriate character output (putchar, console write, etc.)
    \item \textbf{Buffering}: Flush output immediately for real-time interaction
    \item \textbf{Character Encoding}: Treat values as ASCII/UTF-8 characters
\end{itemize}

\section{Porting Guidelines}

\subsection{Platform Requirements}
\textbf{Minimum Requirements}:
\begin{itemize}
    \item 8-bit arithmetic operations
    \item 64KB RAM or equivalent storage mechanism
    \item File I/O capability for loading .mhfb files
    \item Character output mechanism
\end{itemize}

\textbf{Recommended Features}:
\begin{itemize}
    \item Error reporting and debugging output
    \item Register state inspection
    \item Execution step-by-step debugging
\end{itemize}

\subsection{Implementation Strategies}

\subsubsection{Software Emulator}
\textbf{Languages}: C, C++, Rust, Go, Python, JavaScript\\
\textbf{Approach}: Direct state machine implementation\\
\textbf{Performance}: Suitable for development and testing

\begin{lstlisting}[language=C, caption=Minimal Emulator Template]
#include <stdint.h>
#include <stdio.h>

typedef struct {
    uint8_t regs[8];
    uint16_t pc;
    uint8_t memory[65536];
    int running;
} CPU;

void execute_instruction(CPU *cpu) {
    uint8_t opcode = cpu->memory[cpu->pc++];
    
    switch (opcode) {
        case 0x00: break; // NOP
        case 0x01: cpu->running = 0; break; // HALT
        // Add other instructions...
        default:
            printf("Unknown opcode: 0x%02X\n", opcode);
            cpu->running = 0;
    }
}
\end{lstlisting}

\subsubsection{FPGA/Hardware Implementation}
\textbf{Components Required}:
\begin{itemize}
    \item 8×8-bit register file
    \item 16-bit program counter
    \item 64KB RAM interface
    \item Simple ALU (8-bit addition minimum)
    \item Instruction decoder
\end{itemize}

\textbf{Control Logic}:
\begin{itemize}
    \item Fetch: Read instruction from memory[PC]
    \item Decode: Parse opcode and operands
    \item Execute: Perform operation
    \item Update: Increment PC or handle branches
\end{itemize}

\subsubsection{Microcontroller Port}
\textbf{Considerations}:
\begin{itemize}
    \item External RAM may be required for full 64KB space
    \item Flash storage for program loading
    \item Serial/UART output for console I/O
    \item Consider reduced memory variants for resource-constrained devices
\end{itemize}

\subsection{Cross-Platform Considerations}

\subsubsection{Endianness}
\textbf{Standard}: Little-endian for multi-byte values\\
\textbf{Address Encoding}: \texttt{addr16 = (high\_byte << 8) | low\_byte}\\
\textbf{File Format}: All multi-byte fields in binary header use host endianness

\subsubsection{File I/O}
\begin{itemize}
    \item Support binary file reading
    \item Handle file size validation
    \item Graceful error handling for missing/corrupt files
\end{itemize}

\subsubsection{Character Output}
\begin{itemize}
    \item Map to platform-specific output (stdio, serial, display)
    \item Support ASCII character set minimum
    \item Consider UTF-8 support for international text
\end{itemize}

\section{Testing and Validation}

\subsection{Test Programs}
Create test .mhfb files for:
\begin{itemize}
    \item Basic arithmetic operations
    \item Memory load/store operations
    \item Console output functionality
    \item Error condition handling
\end{itemize}

\subsection{Reference Behavior}
All implementations should produce identical output to mhfsoft.c for the same input programs.

\subsection{Debugging Features}
Recommended debugging capabilities:
\begin{itemize}
    \item Register state display
    \item Memory inspection
    \item Step-by-step execution
    \item Instruction trace logging
\end{itemize}

\section{Extensions and Future Development}

\subsection{Instruction Set Extensions}
Reserved opcode ranges for future expansion:
\begin{itemize}
    \item \textbf{0x02-0x0F}: System/control instructions
    \item \textbf{0x12-0x1F}: Memory operations
    \item \textbf{0x21-0x2F}: Arithmetic operations
    \item \textbf{0x30-0x3F}: Logic operations
    \item \textbf{0x40-0x4F}: Branch/jump operations
\end{itemize}

\subsection{Potential Enhancements}
\begin{itemize}
    \item Stack operations and SP register
    \item Conditional branches and status flags
    \item 16-bit arithmetic mode
    \item Interrupt handling
    \item Multiple I/O ports
    \item DMA support
\end{itemize}

\section{Reference Implementation}

The canonical reference implementation is mhfsoft.c, which demonstrates:
\begin{itemize}
    \item MHF binary format loading
    \item Complete instruction execution
    \item Memory-mapped I/O handling
    \item Error handling and reporting
\end{itemize}

All ports should strive for functional compatibility with this reference implementation.

\section{Conclusion}

The mighf system provides a clean, minimal CPU architecture suitable for educational purposes, embedded applications, and experimental computing projects. Its simplicity enables rapid porting to new platforms while maintaining consistent behavior across implementations.

For questions, bug reports, or contributions to the mighf ecosystem, refer to the project documentation and community resources.

\end{document}
